\documentclass[a4paper,12pt]{article}
\usepackage[T2A]{fontenc}
\usepackage[utf8]{inputenc}
\usepackage[russian]{babel}
\usepackage{amsmath,amssymb,mathtools}
\usepackage{icomma}
\usepackage[margin=22mm,top=22mm,bottom=24mm]{geometry}
\usepackage{microtype}
\usepackage{tikz}
\usetikzlibrary{calc,arrows.meta,positioning,shapes.geometric,angles,quotes}
\usepackage{tkz-euclide}
\usepackage{booktabs,tabularx,longtable,array}
\usepackage{enumitem}
\usepackage{hyperref}
\usepackage{parskip}
\usepackage{fancyhdr}
\usepackage{setspace}
\usepackage{xcolor}
\definecolor{accent}{HTML}{008080}
\definecolor{darkaccent}{HTML}{004D4D}
\definecolor{textgray}{HTML}{333333}
\definecolor{softgray}{HTML}{697478}
\definecolor{paleblue}{HTML}{DDEDED}
\hypersetup{colorlinks=true,linkcolor=darkaccent,urlcolor=darkaccent,pdfauthor={Лёвин Артём Александрович},pdftitle={Чек-лист. Кинематика на наклонной плоскости и упругий удар}}
\setlength{\headheight}{14pt}
\setlength{\parindent}{0pt}
\setlength{\parskip}{0.48em}
\linespread{1.06}
\setlist[itemize]{leftmargin=1.4em,itemsep=2pt,topsep=3pt}
\setlist[enumerate]{leftmargin=1.65em,itemsep=4pt,topsep=4pt}
\fancyhf{}
\fancyhead[L]{\small\color{darkaccent} Кинематика и упругий удар}
\fancyhead[R]{\small\color{softgray} София Лузянина}
\fancyfoot[C]{\footnotesize\color{softgray} Лёвин Артём Александрович \textperiodcentered\ эксклюзивно для Софии Лузяниной}
\fancyfoot[R]{\footnotesize\color{softgray}\thepage}
\renewcommand{\headrulewidth}{0.25pt}
\renewcommand{\footrulewidth}{0pt}
\pagestyle{fancy}
\newcommand{\key}[1]{\textcolor{darkaccent}{\textbf{#1}}}
\newcommand{\checkit}{\(\square\)}
\newcommand{\vect}[1]{\vec{#1}}
\newcommand{\gacc}{g\sin\alpha}
\newcommand{\ex}[1]{\vspace{0.28em}\textbf{#1}\quad}
\newcommand{\separate}{\par\vspace{0.25em}\noindent\color{accent}\rule{\linewidth}{0.5pt}\color{black}\par\vspace{0.25em}}
\newcommand{\secstyle}[1]{\section*{\textcolor{darkaccent}{#1}}}
\newcommand{\smallhead}[1]{\subsection*{\textcolor{darkaccent}{#1}}}
\newcommand{\tagline}[1]{\begin{center}\small\color{softgray}#1\end{center}}
\begin{document}

% --- Cover ---
\thispagestyle{empty}
\begin{center}
\vspace*{23mm}
{\large\color{darkaccent}\MakeUppercase{Физика}}\par
\vspace{18mm}
{\LARGE Чек-лист}\par
\vspace{9mm}
{\LARGE\color{accent}\bfseries Кинематика на наклонной плоскости\\[5pt]и абсолютно упругий удар}\par
\vspace{14mm}
{\large Движение шайбы \textperiodcentered\ отражение шарика от плиты}\par
{\normalsize Сочетание кинематики, геометрии и закона сохранения энергии}\par
\vfill
{\large София Лузянина}\par
\vspace{10mm}
{\large 20 ноября 2026 г.}\par
\vspace{13mm}
{\small Лёвин Артём Александрович}\\[4pt]
{\small эксклюзивно для Софии Лузяниной}\par
\vspace{22mm}
\begin{tikzpicture}[x=1cm,y=1cm]
\draw[accent,line width=0.8pt] (0,0) .. controls (1.25,0.65) and (2.3,-0.48) .. (3.5,0)
 .. controls (4.7,0.53) and (5.6,-0.50) .. (6.9,0)
 .. controls (8.0,0.52) and (8.9,-0.34) .. (10.4,0);
\fill[darkaccent] (4.9,0.04) circle (1.2pt);
\end{tikzpicture}
\end{center}
\clearpage

\secstyle{Что нужно уверенно уметь}
\begin{itemize}[label={}]
\item \checkit\ Выбирать оси так, чтобы они отражали реальные независимые направления движения тела.
\item \checkit\ Отличать \emph{координату} от \emph{расстояния}, а \emph{модуль скорости} --- от её проекции на ось.
\item \checkit\ Вычислять компоненты начальной скорости и ускорения; составлять уравнения \(x(t),y(t),v_x(t),v_y(t)\).
\item \checkit\ Находить момент возвращения в исходную координату и момент, когда составляющая скорости вдоль выбранной оси равна нулю.
\item \checkit\ Обосновывать геометрически направления скорости до и после абсолютно упругого удара.
\item \checkit\ Согласованно применять кинематику и закон сохранения механической энергии.
\end{itemize}
\separate
\secstyle{1. Шайба на наклонной плоскости}
\textbf{Ситуация.} Плоскость наклонена к горизонту под углом \(\alpha\) и пересекает горизонтальную плоскость по прямой \(AB\). Шайбу запускают из точки \(A\) по поверхности с начальной скоростью \(v_0\). Скорость составляет угол \(\beta\) с прямой \(AB\). Трением пренебрегаем; требуется определить точку \(B\), в которой шайба вновь окажется на прямой \(AB\).

\begin{center}
\begin{tikzpicture}[scale=0.92,>=Latex]
\tkzDefPoint(0,0){A}
\tkzDefPoint(7.2,0){B}
\tkzDefPoint(1.8,1.04){V}
\tkzDefPoint(0,3.0){C}
\tkzDrawSegment[line width=1.05pt,color=black](A,B)
\draw[softgray,->] (0,0)--(7.95,0) node[right]{$x$};
\draw[softgray,->] (0,0)--(0,3.10) node[above]{$y$ (вверх по склону)};
\draw[accent,line width=1.4pt,->] (A) .. controls (2.10,1.21) and (5.10,1.21) .. (B);
\draw[darkaccent,line width=1.15pt,->] (A)--(V) node[above right]{$\vect v_0$};
\tkzDrawPoints[fill=darkaccent,size=3](A,B)
\tkzLabelPoint[below left](A){$A$}
\tkzLabelPoint[below right](B){$B$}
\draw[darkaccent] (0.9,0) arc (0:30:0.9);
\node[darkaccent] at (1.25,0.34) {$\beta$};
\node[softgray,align=center,font=\small] at (4.4,3.0) {Схема в координатах,\\лежащих \emph{в плоскости} склона};
\end{tikzpicture}
\end{center}
\textbf{Как выбирать оси.} Ось \(Ox\) направим вдоль \(AB\), а ось \(Oy\) --- \emph{внутри наклонной плоскости} вверх по линии наибольшего подъёма. Тогда точка \(A\) имеет координаты \((0,0)\), а точка возвращения \(B\) лежит на \(y=0\).

\key{Почему ось \(Oy\) здесь направлена вдоль склона?} Шайба движется по поверхности сразу в двух направлениях: вдоль \(AB\) и поперёк \(AB\) вверх/вниз по склону. В динамике при разборе сил часто берут вторую ось \emph{перпендикулярно} поверхности; для поиска двух координат движения по поверхности удобнее выбранные здесь оси.

\smallhead{Что делает сила тяжести}
Ускорение свободного падения \(\vect g\) направлено вертикально вниз. Его проекция \emph{на поверхность} направлена вниз по склону и имеет модуль \(g\sin\alpha\). Вдоль \(AB\) ускорения нет. Поэтому
\[
\boxed{a_x=0,\qquad a_y=-g\sin\alpha.}
\]
Перпендикулярная к плоскости составляющая силы тяжести компенсируется нормальной реакцией опоры: \(N=mg\cos\alpha\). Она не влияет на координаты \(x,y\) при отсутствии трения.

\begin{center}
\begin{tikzpicture}[scale=0.87,>=Latex]
\tkzDefPoint(-3,0){O}
\tkzDefPoint(3,2.85){K}
\tkzDrawSegment[line width=1pt](O,K)
\draw[softgray,thin] (-3,0)--(3.7,0);
\draw[darkaccent,->,line width=1.1pt] (-.4,1.25)--(-.4,-1.8) node[below]{$m\vect g$};
\draw[accent,->,line width=1.1pt] (-.4,1.25)--(-2.35,.32) node[above left]{$mg\sin\alpha$};

\draw[darkaccent] (-2.05,0) arc (0:25.5:0.95);
\node[darkaccent] at (-1.48,.25) {$\alpha$};
\node[font=\small,softgray,align=center] at (4.5,.6) {Вид сбоку:\\ускорение вдоль\\склона равно\\$g\sin\alpha$};
\end{tikzpicture}
\end{center}

\smallhead{Четыре уравнения, которые нужны в задаче}
Сначала разложим начальную скорость: \(v_{0x}=v_0\cos\beta\), \(v_{0y}=v_0\sin\beta\). Затем применим формулы равноускоренного движения к каждой оси отдельно:
\[
\begin{aligned}
\text{Ось }Ox:\quad &v_x(t)=v_0\cos\beta,&\quad x(t)&=v_0\cos\beta\,t,\\[3pt]
\text{Ось }Oy:\quad &v_y(t)=v_0\sin\beta-g\sin\alpha\,t,&
 y(t)&=v_0\sin\beta\,t-\frac{g\sin\alpha}{2}t^2.
\end{aligned}
\]
\key{Проверка знаков.} \(v_{0y}>0\), потому что шайба сначала движется вверх по склону. \(a_y<0\), потому что проекция тяжести направлена в противоположную сторону.

\smallhead{Как получить расстояние \(AB\)}
При повторном пересечении прямой \(AB\) поперечная координата снова равна нулю:
\[
0=y(t)=v_0\sin\beta\,t-\frac{g\sin\alpha}{2}t^2
=t\left(v_0\sin\beta-\frac{g\sin\alpha}{2}t\right).
\]
Первый корень \(t=0\) соответствует запуску. Второй --- искомому возвращению:
\[
\boxed{t_{\text{возвр}}=\frac{2v_0\sin\beta}{g\sin\alpha}.}
\]
Вдоль \(Ox\) шайба движется равномерно, поэтому расстояние по прямой \(AB\) равно
\[
\boxed{AB=x(t_{\text{возвр}})=\frac{v_0^2\sin(2\beta)}{g\sin\alpha}.}
\]
Здесь использована формула \(2\sin\beta\cos\beta=\sin(2\beta)\).

\smallhead{Ещё один важный момент: вершина траектории}
На самом высоком положении \emph{вдоль склона} составляющая \(v_y\) обращается в нуль:
\[
 v_y(t_{\text{верх}})=0,\quad
 t_{\text{верх}}=\frac{v_0\sin\beta}{g\sin\alpha},\quad
 y_{\max}=\frac{v_0^2\sin^2\beta}{2g\sin\alpha}.
\]
Это расстояние \(y_{\max}\), измеренное \emph{вдоль поверхности}. Его вертикальная составляющая равна \(\Delta h=y_{\max}\sin\alpha\). На вершине \(v_y=0\), однако \(v_x=v_0\cos\beta\) сохраняется, поэтому полная скорость обычно \emph{не} равна нулю.

\textbf{Условия.} Для описанного возвращения берём \(0<\alpha<90^\circ\), \(0<\beta<90^\circ\), \(v_0>0\); бесконечная плоскость предполагается достаточно большой. При \(\alpha=0\) сила тяжести не возвращает шайбу на \(AB\).

\smallhead{Разобранный пример}
Шайба запущена с \(v_0=10\,\text{м/с}\), \(\beta=30^\circ\); наклон плоскости \(\alpha=30^\circ\), \(g=10\,\text{м/с}^2\). Найдём время движения и расстояние \(AB\):
\[
 t_{\text{возвр}}=\frac{2\cdot10\cdot\sin30^\circ}{10\sin30^\circ}=2\,\text{с},\qquad
 AB=10\cos30^\circ\cdot2=10\sqrt3\,\text{м}.
\]
\key{Смысл проверки:} время положительно; оба момента \(t=0\) и \(t=2\) дают \(y=0\); единицы дальности --- метры.

\clearpage
\thispagestyle{fancy}
\begin{center}
{\Large\color{darkaccent}\textbf{Алгоритм: найти возвращение шайбы на AB}}
\vspace{11mm}
\end{center}
\tikzset{
 st/.style={draw=darkaccent,rounded corners=8pt,minimum height=11mm,text width=92mm,align=center,fill=paleblue!30,font=\normalsize},
 op/.style={draw=softgray,rounded corners=2pt,minimum height=13mm,text width=108mm,align=center,font=\normalsize},
 dec/.style={draw=darkaccent,diamond,aspect=2.6,inner xsep=2pt,inner ysep=3pt,text width=52mm,align=center,font=\small},
 arr/.style={-Latex,line width=.95pt,color=darkaccent}}
\begin{center}
\begin{tikzpicture}[node distance=12mm and 21mm]
\node[st] (start) {Начало: известны \(v_0,\alpha,\beta\)};
\node[op,below=of start] (axis) {Выбираем \(Ox\) вдоль \(AB\), \(Oy\) вверх по склону; записываем \(a_x=0\), \(a_y=-g\sin\alpha\)};
\node[op,below=of axis] (eq) {Записываем \(x(t)=v_0\cos\beta\,t\) и \(y(t)=v_0\sin\beta\,t-\tfrac12g\sin\alpha\,t^2\)};
\node[op,below=of eq] (root) {Для точки на \(AB\) полагаем \(y(t)=0\). Получаем корни \(t=0\) и \(t=\dfrac{2v_0\sin\beta}{g\sin\alpha}\)};
\node[dec,below=15mm of root] (judge) {Корень \(t>0\)?};
\node[op,right=22mm of judge,text width=35mm,font=\small] (discard) {\(t=0\): момент старта; не используем как время возвращения};
\node[op,below=19mm of judge] (dist) {Подставляем положительный корень в \(x(t)\): \(AB=\dfrac{v_0^2\sin2\beta}{g\sin\alpha}\)};
\node[st,below=of dist] (end) {Ответ: время возвращения и расстояние \(AB\)};
\draw[arr] (start)--(axis);
\draw[arr] (axis)--(eq);
\draw[arr] (eq)--(root);
\draw[arr] (root)--(judge);
\draw[arr] (judge)--node[right,font=\small] {да} (dist);
\draw[arr] (judge)--node[above,font=\small] {нет} (discard);
\draw[arr] (discard.north) |- ($(root.east)+(0,0)$);
\draw[arr] (dist)--(end);
\end{tikzpicture}
\end{center}
\clearpage

\secstyle{2. Упругий отскок шарика от наклонной плиты}
\textbf{Ситуация.} Шарик отпущен без начальной скорости с высоты \(H\) над горизонтальной землёй. Через время \(t\) он попадает на неподвижную плиту, наклонённую под углом \(\alpha\) к горизонту. Удар считаем абсолютно упругим и гладким, сопротивлением воздуха пренебрегаем. Требуется определить наибольшую высоту \(h\) над землёй после удара.

\smallhead{Шаг 1. Движение до удара}
Ось высот направим \emph{вертикально вверх}. Тогда за время \(t\) шарик прошёл вниз \(s=gt^2/2\), его высота в момент столкновения и модуль скорости равны
\[
\boxed{h_0=H-\frac{gt^2}{2}},\qquad \boxed{v=gt}.
\]
В первом равенстве \(h_0\) --- \emph{высота над землёй}; \(gt^2/2\) --- \emph{положительное расстояние падения}. Если пользоваться \emph{координатой} \(y\), её изменение при движении вниз отрицательно.

\smallhead{Шаг 2. Сначала обоснуем угол отражения}
Нормаль --- прямая, перпендикулярная плите. Поскольку плита образует с горизонталью угол \(\alpha\), нормаль образует с вертикалью \emph{тот же угол} \(\alpha\) (угол между двумя прямыми равен углу между их перпендикулярами).

При идеальном гладком упругом ударе нормальная составляющая скорости меняет знак, касательная сохраняется. Поэтому угол падения и угол отражения \emph{относительно нормали} совпадают. Здесь оба равны \(\alpha\). Следовательно, направление скорости после удара образует с вертикалью угол \(2\alpha\).

\begin{center}
\begin{tikzpicture}[scale=1.02,>=Latex]
\tkzDefPoint(-3,1.75){L}
\tkzDefPoint(3,-1.75){R}
\tkzDefPoint(0,0){O}
\tkzDrawSegment[line width=1.25pt](L,R)
\draw[softgray] (-3,0)--(3.45,0) node[right]{горизонт};
\draw[softgray,dashed,thin] (O)--(0,3.3) node[above]{вертикаль};
\draw[softgray,dashed,thin] (O)--(0,-.65);
\draw[darkaccent,->,line width=1.3pt] (0,3.0)--(0,0.15) node[midway,left]{$\vect v_{\text{до}}$};
\draw[darkaccent,->,line width=1.2pt] (O)--(1.5,2.55) node[above right]{нормаль};
\draw[accent,->,line width=1.65pt] (O)--(2.9,1.67) node[above right]{$\vect v_{\text{после}}$};
\draw[darkaccent] (-1.12,0) arc (180:150:1.12);
\node[darkaccent] at (-1.5,.23) {$\alpha$};
\draw[accent] (0,1.15) arc (90:30:1.15);
\node[accent] at (.76,1.08) {$2\alpha$};
\node[softgray,font=\small] at (2.35,-1.6) {плита};
\tkzDrawPoints[fill=darkaccent,size=3](O)
\tkzLabelPoint[below left](O){$O$}
\end{tikzpicture}
\end{center}
На рисунке плита наклонена вниз вправо, поэтому шарик уходит вправо. При зеркальном повороте рисунка направление горизонтальной составляющей поменяется, но \emph{максимальная высота останется той же}.

Если \(0<\alpha<45^\circ\), после столкновения шарик действительно движется вверх, а угол вылета к горизонту равен
\[
\boxed{\varphi=90^\circ-2\alpha.}
\]
\key{Не перепутайте углы:} \(\alpha\) измеряется от нормали для падающего луча; \(2\alpha\) --- от вертикали для отражённой скорости; \(\varphi\) --- от горизонта.

\smallhead{Шаг 3. Проекции скорости после удара}
При упругом ударе о \emph{неподвижную гладкую} плиту модуль скорости сохраняется: \(v'=v=gt\). Выбрав ось \(Ox\) в направлении горизонтального движения после удара, получаем
\[
\boxed{v'_x=gt\sin(2\alpha)},\qquad
\boxed{v'_y=gt\cos(2\alpha)}.
\]
После отскока это обычное движение тела под углом к горизонту с ускорениями \(a_x=0\), \(a_y=-g\):
\[
 x(\tau)=v'_x\tau,\quad
 y(\tau)=h_0+v'_y\tau-\frac{g\tau^2}{2}.
\]
Здесь \(\tau\) --- время, отсчитываемое \emph{от момента удара}; оно отличается от времени свободного падения \(t\).

\smallhead{Шаг 4. Высота в вершине траектории}
Для \(0<\alpha<45^\circ\) шарик поднимается до момента \(v_y(\tau)=v'_y-g\tau=0\). Значит,
\[
\tau_{\text{верх}}=t\cos(2\alpha),\qquad
\Delta h=\frac{(v'_y)^2}{2g}=\frac{gt^2\cos^2(2\alpha)}{2}.
\]
Прибавим подъём после удара к высоте места удара:
\[
\begin{aligned}
 h&=h_0+\Delta h
 =H-\frac{gt^2}{2}+\frac{gt^2\cos^2(2\alpha)}{2}\\
 &=\boxed{H-\frac{gt^2}{2}\sin^2(2\alpha)}.
\end{aligned}
\]

\smallhead{Тот же результат через закон сохранения энергии}
Высшая точка после отскока --- это точка, где \(v_y=0\), но \(v_x=v'_x\ne0\). Значит, \emph{часть кинетической энергии остаётся}:
\[
mgH=mgh+\frac{m(v'_x)^2}{2}
=mgh+\frac{m g^2t^2\sin^2(2\alpha)}{2}.
\]
Сокращая на \(mg\), снова получаем формулу для \(h\). Писать здесь \(mgH=mgh\) нельзя: это означало бы ошибочно считать, что в вершине остановилось всё тело.

\smallhead{Проверка физического смысла и крайних случаев}
\begin{itemize}
\item \(h_0\ge0\), иначе столкновение произошло бы ниже уровня земли: нужно \(0<t\le\sqrt{2H/g}\).
\item Если \(\alpha=0\), идеально горизонтальная плита возвращает шарик к исходной высоте: \(h=H\).
\item Если \(\alpha=45^\circ\), отражённая скорость горизонтальна, высота после удара сразу максимальна: \(h=h_0\).
\item Если \(45^\circ<\alpha<90^\circ\), \(v'_y<0\): после удара шарик уже движется вниз. В таком случае \emph{максимальная высота после удара} равна \(h_0\); формула с \(\sin^2(2\alpha)\) вычисляет вершину продолженной параболы при отрицательном времени и для этой задачи неприменима.
\end{itemize}

Для диапазона \(0\le\alpha<90^\circ\) все случаи объединяются так:
\[
\boxed{h_{\max}=H-\frac{gt^2}{2}+\frac{gt^2}{2}\bigl[\max\{0,\cos(2\alpha)\}\bigr]^2.}
\]

\smallhead{Разобранный пример}
Пусть \(H=40\,\text{м}\), \(t=2\,\text{с}\), \(\alpha=30^\circ\), \(g=10\,\text{м/с}^2\). До удара шарик опустился на \(20\) м: \(h_0=20\) м, \(v=20\) м/с. После удара \(v'_y=20\cos60^\circ=10\) м/с. Дальнейший подъём составляет \(\Delta h=10^2/(2\cdot10)=5\) м, следовательно,
\[
\boxed{h=20+5=25\,\text{м}.}
\]
Энергетическая проверка даёт то же: \(40-25=15=(20\sin60^\circ)^2/(2\cdot10)\).

\clearpage
\secstyle{3. Типичные ошибки и как их обнаружить}
\begin{enumerate}[label=\textcolor{darkaccent}{\arabic*.}]
\item \textbf{Оси не подписаны.} Прежде чем ставить плюс или минус у ускорения и скорости, обязательно обозначьте начало координат и положительные направления \(Ox,Oy\).
\item \textbf{В задаче о шайбе выбирают нормаль к плоскости вместо второй оси движения.} Для двух координат на поверхности используйте направления \(AB\) и подъём по склону.
\item \textbf{Путают два нуля:} \(y=0\) определяет возвращение на \(AB\), а \(v_y=0\) --- вершину траектории. Во второй точке \(v_x\) сохраняется.
\item \textbf{Выбирают \(t=0\) как время полёта.} Корень \(0\) означает старт; физически интересен второй положительный корень.
\item \textbf{Сразу пишут \(2\alpha\) после отскока.} Сначала покажите нормаль и докажите, что она составляет с вертикалью угол \(\alpha\), затем используйте равенство углов относительно нормали.
\item \textbf{Приравнивают высоту и отрицательное перемещение.} Величина высоты над землёй неотрицательна, тогда как изменение вертикальной \emph{координаты} при падении отрицательно.
\item \textbf{В энергетическом уравнении берут \(v_x\) как полную скорость в момент столкновения.} Во время падения \(v=gt\), а в вершине после отскока \(v_x=gt\sin2\alpha\). Это разные моменты движения.
\item \textbf{Применяют формулу подъёма при \(\alpha>45^\circ\).} Сначала проверьте знак \(v'_y=gt\cos2\alpha\): для отрицательного значения дальнейшего подъёма нет.
\end{enumerate}

\smallhead{Короткая памятка о разных «нулевых условиях»}
\renewcommand{\arraystretch}{1.35}
\begin{tabularx}{\linewidth}{@{}p{.28\linewidth}X@{}}
\toprule
\textbf{Что обнуляется} & \textbf{Что это означает физически} \\
\midrule
\(y(t)=0\) у шайбы & Возвращение на прямую \(AB\) (после старта).\\
\(v_y(t)=0\) у шайбы & Максимальное смещение вверх по склону.\\
\(v_y(\tau)=0\) у шарика & Верхняя точка траектории после отскока, если шарик летел вверх.\\
\(v_x=0\) после удара & Только для специальных углов; общее условие вершины этому не равно.\\
\bottomrule
\end{tabularx}
\separate
\key{Формат решения повышенной сложности:}\par\begin{center}\small
схема с осями $\longrightarrow$ геометрия углов $\longrightarrow$ проекции\\
$\longrightarrow$ уравнения $\longrightarrow$ физический выбор корня или экстремума\\
$\longrightarrow$ проверка единиц и условий.
\end{center}

\clearpage
\secstyle{Тренировочный блок. Путь А --- 10 задач}
\tagline{Решайте последовательно. В каждом задании сначала нарисуйте схему и подпишите выбранные оси.}
\textbf{Условия для задач 1--5.} Наклонная плоскость образует с горизонтальной угол \(\alpha\), линия её пересечения с горизонтальной плоскостью --- \(AB\). Шайба начинает движение из \(A\) со скоростью \(v_0\) под углом \(\beta\) к \(AB\), сначала поднимаясь по склону. Трения нет, \(g=10\,\text{м/с}^2\).

\begin{enumerate}[label=\textbf{\arabic*.},itemsep=10pt]
\item \textbf{Проекции и законы движения.} Шайба стартует со скоростью \(8\,\text{м/с}\) под углом \(30^\circ\) к прямой \(AB\). Угол наклона поверхности к горизонту также равен \(30^\circ\). Выберите ось \(Ox\) вдоль \(AB\), а ось \(Oy\) вверх по поверхности. Найдите численные выражения для \(v_{0x},v_{0y},a_x,a_y\) и запишите уравнения \(x(t),y(t)\), где координаты отсчитываются от точки старта.
\item \textbf{Возвращение на линию \(AB\).} Шайбу запустили со скоростью \(12\,\text{м/с}\) под углом \(45^\circ\) к \(AB\); угол наклона поверхности составляет \(30^\circ\). Вычислите время от запуска до повторного пересечения \(AB\), затем найдите расстояние от \(A\) до точки возвращения \(B\). Поясните, почему корень \(t=0\) из уравнения \(y(t)=0\) не подходит в качестве времени движения.
\item \textbf{Наивысшее положение на склоне.} Шайба стартует со скоростью \(14\,\text{м/с}\) под углом \(30^\circ\) к \(AB\) на плоскости с \(\alpha=30^\circ\). Определите время до верхней точки, максимальную \emph{координату вдоль склона} \(y_{\max}\) и прирост \emph{вертикальной} высоты над уровнем линии \(AB\). Объясните, почему эти два расстояния различаются.
\item \textbf{Восстановление угла запуска.} Скорость шайбы равна \(10\,\text{м/с}\), наклон поверхности составляет \(30^\circ\), а измеренное расстояние \(AB\) равно \(20\,\text{м}\). Найдите угол \(\beta\) между начальной скоростью и линией \(AB\), предполагая \(0^\circ<\beta<90^\circ\). Обязательно запишите уравнение, из которого определяется угол.
\item \textbf{Сравнение двух запусков.} С одной точки наклонной плоскости с \(\alpha=30^\circ\) две одинаковые шайбы запускают с одной скоростью \(10\,\text{м/с}\), но под углами \(\beta_1=30^\circ\) и \(\beta_2=60^\circ\) к \(AB\). Для каждой шайбы найдите время возвращения на \(AB\) и расстояние до точки возвращения. Объясните, как могут совпадать расстояния при разных временах движения.
\end{enumerate}

\clearpage
\textbf{Условия для задач 6--10.} Шарик отпускают без начальной скорости с высоты \(H\) над горизонтальной землёй. Через время \(t\) он сталкивается с неподвижной гладкой плитой, составляющей угол \(\alpha\) с горизонтом. Удар идеально упругий; сопротивлением воздуха пренебречь, \(g=10\,\text{м/с}^2\). Плита ориентирована как на рисунке в теоретическом разделе (падает вертикально вниз, после удара при \(\alpha<45^\circ\) летит вправо и вверх).

\begin{enumerate}[label=\textbf{\arabic*.},start=6,itemsep=11pt]
\item \textbf{Скорость и направление после отражения.} Шарик падает с высоты \(35\,\text{м}\) и ударяется о плиту через \(2\,\text{с}\); \(\alpha=30^\circ\). Найдите высоту точки удара, модуль скорости непосредственно перед и после столкновения, горизонтальную и вертикальную составляющие скорости после удара. Из геометрии нормали определите угол вылета к горизонту. Обоснуйте выбор знаков компонент.
\item \textbf{Максимальная высота после удара.} Из высоты \(H=20\,\text{м}\) шарик падает в течение \(1\,\text{с}\) и затем отражается от плиты с \(\alpha=30^\circ\). Сначала найдите высоту места удара; затем вычислите вертикальную составляющую скорости после удара, добавочную высоту подъёма и итоговую максимальную высоту над землёй. Проверьте результат формулой закона сохранения энергии.
\item \textbf{Другой наклон.} Шарик падает из точки на высоте \(45\,\text{м}\) и спустя \(2\,\text{с}\) сталкивается с плитой, наклонённой к горизонту на \(15^\circ\). Постройте схему отражения с нормалью, определите угол вылета к горизонту и найдите наибольшую высоту после удара. Поясните, почему для вычисления подъёма используется \(\cos(2\alpha)\).
\item \textbf{Восстановление угла плиты.} Шарик падает с высоты \(H=30\,\text{м}\); через \(2\,\text{с}\) происходит упругий удар о наклонную плиту. Экспериментально установлено, что после удара максимальная высота шарика равна \(25\,\text{м}\). Найдите угол наклона \(\alpha\), если заранее известно, что \(0^\circ<\alpha<45^\circ\). Проверьте полученный угол подстановкой в формулу высоты.
\item \textbf{Проверка условий применимости.} Два шарика по отдельности падают с высоты \(H=50\,\text{м}\) и ударяются о две разные гладкие плиты через \(2\,\text{с}\). Первая плита наклонена на \(30^\circ\), вторая --- на \(50^\circ\). Для каждого случая найдите вертикальную составляющую скорости \emph{сразу после} удара и наибольшую высоту над землёй, достигнутую после удара (включая момент столкновения). Объясните, почему одну и ту же формулу подъёма нельзя безоговорочно использовать для обоих углов.
\end{enumerate}

\clearpage
\secstyle{Ответы к тренировочному блоку}
\tagline{Сверяйте после самостоятельного решения. В заданиях с несколькими величинами порядок ответов совпадает с порядком вопросов.}
\renewcommand{\arraystretch}{1.45}
\small
\begin{longtable}{@{}p{.065\linewidth}p{.87\linewidth}@{}}
\toprule
\textbf{№} & \textbf{Ответ}\\
\midrule
\endfirsthead
\toprule
\textbf{№} & \textbf{Ответ (продолжение)}\\
\midrule
\endhead
1 & \(v_{0x}=4\sqrt3\,\text{м/с}\), \(v_{0y}=4\,\text{м/с}\), \(a_x=0\), \(a_y=-5\,\text{м/с}^2\); \(x(t)=4\sqrt3\,t\), \(y(t)=4t-\tfrac52t^2\) (СИ).\\
2 & \(t=\dfrac{12\sqrt2}{5}\,\text{с}\approx3{,}39\,\text{с}\); \(AB=\dfrac{144}{5}\,\text{м}=28{,}8\,\text{м}\); \(t=0\) --- старт.\\
3 & \(t_{\text{верх}}=1{,}4\,\text{с}\); \(y_{\max}=4{,}9\,\text{м}\) вдоль склона; \(\Delta h=2{,}45\,\text{м}\) вертикально.\\
4 & \(\sin(2\beta)=1\), поэтому \(\beta=45^\circ\).\\
5 & \(t_1=2\,\text{с}\), \(t_2=2\sqrt3\,\text{с}\); \(AB_1=AB_2=10\sqrt3\,\text{м}\). Углы \(30^\circ\) и \(60^\circ\) дают одинаковый \(\sin(2\beta)\).\\
6 & \(h_0=15\,\text{м}\); \(v_{\text{до}}=v_{\text{после}}=20\,\text{м/с}\); \(v'_x=10\sqrt3\,\text{м/с}\), \(v'_y=10\,\text{м/с}\); \(\varphi=30^\circ\).\\
7 & \(h_0=15\,\text{м}\), \(v'_y=5\,\text{м/с}\), \(\Delta h=1{,}25\,\text{м}\), \(h=16{,}25\,\text{м}\).\\
8 & \(\varphi=60^\circ\); \(h=40\,\text{м}\).\\
9 & \(\sin^2(2\alpha)=\tfrac14\); с учётом \(0^\circ<\alpha<45^\circ\): \(\alpha=15^\circ\).\\
10 & \(v'_{y,1}=10\,\text{м/с}\), \(h_1=35\,\text{м}\); \(v'_{y,2}=20\cos100^\circ\,\text{м/с}<0\), \(h_2=30\,\text{м}\) (непосредственно при ударе).\\
\bottomrule
\end{longtable}
\normalsize
\separate
\begin{center}
\textcolor{darkaccent}{\textbf{Итоговая связка методов}}\\[3pt]
геометрия углов $\longrightarrow$ проекции скорости $\longrightarrow$ кинематика $\longrightarrow$ физический смысл ответа.
\end{center}
\end{document}
